\section{Discussion}

Subsections 7.1 to 7.5 are written mainly from the exploratory results of Section~6.1, which come from single runs on the validation sets. The confirmatory comparison HF-c supports the stop removal of B2 on the sealed test set (Section~6.2), and the statements below about B1 and B3 remain exploratory. The limitations are in Subsection 7.6. The section does not repeat the numerical results.

\subsection{Why the filter reduces protective stops}
The certified layer triggers when a valid scan return lies inside the protective field of the speed regime the robot is in. The filter constrains each command so that, over one prediction horizon, no valid return lies inside the field of the regime the robot will be in. A command that would trigger the layer is therefore replaced before the layer sees it, and the layer's own stop is not needed. The results are consistent with this mechanism. With the filter, the protective stops of B1, B2 and B3 fall from 638, 591 and 78 to 0, 4 and 0 on routes without an unmapped obstacle, and from 624, 638 and 456 to 12, 16 and 0 on routes with one (Section~6.1). The share of steps with a latched stop falls to 0.002 or less, and every unrecoverable stop disappears. Across the three controllers the filter changes the command on between 19\,\% and 89\,\% of the steps, but it finds no feasible regime, and so has to brake at the maximum rate, on at most 0.4\,\% of them. Most corrections are therefore ordinary changes of the command and not emergency braking.

These outcomes show that the stops fall, and they do not by themselves show why. The remaining stops, 12 of B1 and 16 of B2 on unmapped routes, are described in Section~6.5. Of the 32 stops that remain across the three controllers, 31 fall in episodes in which the filter fell back to maximum braking at least once. Whether the speed- and direction-dependent geometry of the field is what removes the stops, rather than any barrier margin on the scan, is the question of the constant-margin comparator of Section~4.4 \tbd{constant-margin comparator}.

\subsection{Where it helps most}
The benefit follows the places where a controller stops. Table~\ref{tab:classical-mapped} and the per-stratum counts of the validation sets show two patterns, both from single runs of 30 routes per cell.

First, the condition matters for success but not for stops. The filter removes the stops both on routes without an unmapped obstacle and on routes with one, but success rises substantially only on the first. On mapped-only routes, success rises from 0 to 90 of 90 for B1 and from 34 to 68 for B2, and on unmapped routes it stays at 6, 38 and 2 of 90 for B3, B1 and B2, because most of these routes end in timeouts (Section~7.3).

Second, the aisle stratum matters for the controller whose stops depend on it. The stops of B3 alone on mapped-only routes fall with the aisle width, from 18.98 per 100~m in narrow aisles to 4.42 in medium and 1.26 in wide aisles, and the filter removes them in every stratum. For B3 the success of narrow aisles rises from 25 to 30, and the other strata were already at 30. B1 and B2 stop at 58 to 71 per 100~m in all three strata, and the filter removes them in all three. Success on mapped-only routes rises in every stratum for B2 (from 9 to 21 of 30 in narrow, from 12 to 26 in medium and from 13 to 21 in wide aisles). These patterns are descriptive. With 30 routes per cell and one run, we give no interval by stratum, and we do not claim that the filter helps more in one stratum than in another beyond the places where the controllers stop.

\subsection{Conservatism and excessive braking}
The filter removes the protective stops, but it does not remove the failure to complete the route. On routes with an unmapped obstacle, the share of timeouts rises when the filter is added: from 0 to 52 of 90 for B1, from 4 to 88 for B2 and from 75 to 84 for B3 (Section~6.1). The unrecoverable stops that ended these episodes alone become timeouts with the filter, so a stop is replaced by a robot that is still free to move but does not reach the goal. This is consistent with a controller that does not replan: the stack follows its route, the obstacle is visible in the scan but absent from the map, and the filter keeps the commanded motion out of the field instead of letting it enter. The filter alters the commands, and the controller has no way to use that information to route around the obstacle. The timeouts of the filtered controllers are described in Section~6.5. None of the 246 timeouts ends below 10\,\% of the planned route, and 149 end between 50 and 90\,\%. The analysis of stalls of barrier filters at undesired equilibria~\cite{mestres2026} suggests that such behaviour is a known property of this kind of filter, but it is a continuous-time result and applies to our discrete, latched setting by analogy only.

On routes without an unmapped obstacle, the filter also lengthens the timeouts of one controller: B2 times out on 22 of 90 routes with it against 10 alone, although its success doubles from 34 to 68. The filter is therefore conservative in some situations, and a reduction in stops must be read beside success and timeouts, as the tables do. The moving speed shows the same caution: B1 moves more slowly with the filter, at 0.128 against 0.180~m/s on mapped-only routes and at 0.053 against 0.149~m/s on unmapped routes, and B2 moves more slowly on unmapped routes (0.032 against 0.058~m/s), whereas B3 changes little (Table~\ref{tab:classical-time}). The comparison is confounded, because the alone versions spend most of their time latched and the filtered versions complete more routes. It indicates that part of the gain is bought with a slower approach, which the stop metric does not show.

\subsection{Dependence on the underlying controller}
The effect of the filter is largest for the controllers whose commands take no account of the protective field. B1, B2 and B3 are field-blind, and the filter removes nearly all of their stops, on routes with and without an unmapped obstacle. B3 is the case that carries the evidence: it is the only one of the three that completes most routes alone (85 of 90 mapped-only routes), so its stops are not a symptom of a controller that cannot navigate. Its stops fall from 8.35 to 0.0 per 100~m on mapped-only routes and from 45.7 to 0.0 on unmapped routes, with no loss of success or moving speed on mapped-only routes.

The two cases with the largest gains are weak baselines. B1 alone has an unrecoverable stop on every one of its 90 mapped routes and completes none, and B2 alone completes 34. The rise of B1 from 0 to 90 successes is therefore not evidence about strong planners. For B3 the filter removes the stops on mapped routes with a change in success of the size of run-to-run variation, and on unmapped routes with no improvement in the ability to complete them. These are single exploratory runs on the validation sets with controllers tuned on the tuning sets, so they are exploratory. The confirmatory comparison on the sealed test set confirms the stop removal for B2 only (Section~6.2), and it does not extend to B1 or B3.

\subsection{Fewer stops versus better navigation}
A reduction in protective stops is not an improvement in navigation by itself, and the results show both cases. On routes without an unmapped obstacle, the fall in stops comes with equal or higher success for B1, B2 and B3 (Section~6.1), and the mission times of B2 and B3 on the routes both versions complete are not longer. There the filter improves the stop rate without a cost that the other measures reveal, with the exception of B2, whose timeouts rise from 10 to 22 of 90 routes while its success doubles. On routes with an unmapped obstacle the picture is different. The stops of B3, B1 and B2 fall to 0.0, 1.2 and 1.6 per 100~m, but the timeouts rise to 84, 52 and 88 of 90, and the success counts stay at 6, 38 and 2. B3 shows the point most clearly: with the filter it has no stop at all, and it completes 6 of 90 routes. A stop rate of zero can be reached by a robot that rarely reaches the goal.

Whether replacing an unrecoverable stop by a timeout is an improvement depends on the operational cost of each. An unrecoverable stop is a robot that cannot legally move again, and a timeout is a robot that can still move but has not arrived. The metrics of this paper count both as a failure to complete the route, and we do not claim that the timeout is the cheaper failure. We therefore restrict the claim to stop removal. The filter reduces protective stops on every controller that stops, it does not by itself make a controller complete routes that require avoiding an obstacle missing from its map, and it can lower the moving speed (Section~7.3). The stop rate must be read beside the success rate, the timeouts and the mission time, as the tables do, and a comparison of controllers on the stop rate alone would favour a controller that does not move. The confirmatory comparison reads the same way: in HF-c the stops of B2 fall from 68.1 to 1.11 per 100~m of planned route, while 147 of 150 routes with the filter end in a timeout.

\subsection{Limitations}
Several limitations should be considered when interpreting the results. First, the PFB operates on a modeled representation of the protective-field geometry and switching logic, which may not capture all implementation-specific characteristics of an industrial safety scanner. Second, the filtering condition is evaluated over one prediction horizon and does not establish continuous-time safety or a formal collision-avoidance guarantee, and the barrier is a candidate whose validity is not shown. Third, the effectiveness of the filter depends on the accuracy of the robot motion model and on the validity of the available range measurements. Fourth, the experiments are in simulation, with a static world and a single platform, and reducing modeled protective stops does not imply equivalent behaviour in a physical certified safety system. Fifth, the evidence is exploratory: the results are single runs on the validation sets with controllers tuned on the tuning sets, the confirmatory comparison HF-c covers B2 alone, and the constant-margin comparator has not been run \tbd{update when the comparator exists}. Sixth, B1 and B2 are weak baselines, so the evidence on a strong controller rests on B3. The proposed method is intended as a navigation-level filtering mechanism, not as a replacement for certified safety functions. \tbd{implications of transferring the approach to a physical AMR}
